% Gaussian mutual information: generated by make_pgfplots.py
\begin{figure*}[t]
\centering
\setlength{\rgplotwidth}{0.245\textwidth}
\begin{tikzpicture}
\begin{axis}[rglegendaxis, legend columns=5]
\addlegendimage{rgdual}\addlegendentry{RG/dual}
\addlegendimage{rgtopk}\addlegendentry{RG/top-$k$}
\addlegendimage{rgopt}\addlegendentry{RG/OPT}
\addlegendimage{rglp}\addlegendentry{RG/LP (B2)}
\addlegendimage{rggamma}\addlegendentry{RG/$\gamma_1$ (B3)}
\addlegendimage{rgmutwo}\addlegendentry{RG/$\mu_2$ (B4)}
\addlegendimage{rgmuthree}\addlegendentry{RG/$\mu_3$ (B4)}
\addlegendimage{rgonee}\addlegendentry{$1/e$}
\end{axis}
\end{tikzpicture}\\[1pt]
\input{\rgroot/plots/gaussian_mi__synthetic_gp_n20_baselines.tex}\hfill
\input{\rgroot/plots/gaussian_mi__synthetic_gp_n100_baselines.tex}\hfill
\input{\rgroot/plots/gaussian_mi__intel_baselines.tex}
\caption{Gaussian mutual information: certified approximation ratio of random greedy (RG, mean over trials) with respect to each upper bound, as a function of the cardinality $k$. $f(S)=I(X_S;X_{V\setminus S})$ for a Gaussian process on random points and for the Intel Berkeley lab temperature sensors. RG/OPT uses the exact optimum (brute force) where it was computable. Additional baselines: the LP over the dual's base sets (B2), the pruned double-greedy bound $\gamma_1$ (B3) and the modular bounds $\mu_2$, $\mu_3$ on the pruned lattice (B4); B3 and B4 bound the unconstrained optimum. Curves are means over graph/sample seeds.}
\label{fig:rg-gaussian_mi-wide-baselines}
\end{figure*}
